\documentclass{article}
\input{../../lib.sty}

\title{Gaussian Differential Privacy Curve}
\author{\AishwaryaRamasethuAuthor, \YuJuKuAuthor}
\begin{document}
\maketitle

\section*{Definition}
For $\mu \geq 0$, the Gaussian differential privacy tradeoff curve is
\[
  \beta_\mu(\alpha) = \Phi\left(\Phi^{-1}(1-\alpha)-\mu\right).
\]
Its privacy profile is
\[
  \delta_\mu(\epsilon) =
  \Phi\left(-\frac{\epsilon}{\mu}+\frac{\mu}{2}\right)
  -e^\epsilon\Phi\left(-\frac{\epsilon}{\mu}-\frac{\mu}{2}\right),
\]
with the value at $\mu=0$ defined by continuity as zero.

\section*{Numerical assumptions}
The implementation is gated by \texttt{idealized-numerics}. Basic interval
arithmetic uses the software directed backend. Native special-function results
remain idealized approximations, widened outward by one native floating-point
unit; this is not a certification of those special functions. Their decreasing
arguments are rounded in the opposite direction before native evaluation.

\section*{Composition}
Independent Gaussian privacy losses compose by addition of their squared
parameters. Thus, for curves with parameters $\mu_i$,
\[
  \mu_{\mathrm{comp}} = \sqrt{\sum_i \mu_i^2}.
\]
The implementation encloses each arithmetic operation with directed interval
arithmetic and rejects a non-finite resulting parameter.

\end{document}
